\chapter{開発環境設定}
\label{ch:setup}

この章では、本格的なLLM実装に先立って必要な開発環境を設定します。 Python、PyTorch、そしてこの本のコードストア構造を見て、最初の「Hello、LLM」プログラムを実行してみましょう。

\section{必要なツール}

\begin{center}
\begin{tabularx}{\textwidth}{lXl}
\toprule
\textbf{ツール}&\textbf{用途}&\textbf{推奨バージョン}\\
\midrule
Python＆実装言語＆3.10以上\\
PyTorch＆ディープラーニングフレームワーク＆2.0以上\\
NumPy＆数値計算＆1.24以上\\
pytest＆ユニットテスト＆7.0以上\\
git＆バージョン管理＆最新\\
\bottomrule
\end{tabularx}
\end{center}

GPUはオプションです。この本のすべてのコードはCPUで動作するように書かれています。ただし、大規模な学習のためにはCUDA対応GPUが強く推奨される。 2巻で分散学習を扱うときは、複数のGPU環境を想定する。

\section{PyTorch インストール}

PyTorchは公式ホームページのインストールガイドに従います。 CPU専用バージョンとCUDAバージョンのうち、自分の環境に合うものを選択する。

\begin{lstlisting}[style=bashstyle]
# CPU のみ（本のすべての例は CPUとして動作）
pip install torch --index-url https://download.pytorch.org/whl/cpu

# CUDA 12.1 サポート(GPU 使用時）
pip install torch --index-url https://download.pytorch.org/whl/cu121
\end{lstlisting}

設置確認は次のようにする。

\begin{lstlisting}[language=Python]
import torch
print(torch.__version__)         # 2.x.x
print(torch.cuda.is_available()) # GPU 利用可能かどうか
\end{lstlisting}

\section{ストレージ構造}

この本のコードは次のような構造を持っています。各章で扱うモジュールは別々のファイルに分かれており、必要な部分だけを読み込むことができます。

\begin{lstlisting}[style=bashstyle]
make-llm-basic/
  src/makellm/
    __init__.py
    tokenizer/              # 第3章：トークナイザー
      base.py
      bpe.py
      unigram.py
    model/                  # 第4〜6章：モデル構造
      embedding.py
      attention.py
      transformer_block.py
      gpt.py
    training/               # 第7章：学習パイプライン
      dataset.py
      loss.py
      optimizers.py
      trainer.py
    inference/              # 第8章：推論エンジン
      sampler.py
      generate.py
    utils/                  # 共通ユーティリティ
      config.py
      seed.py
      logging.py
  tests/                    # pytest ユニットテスト
  book/                     # この本の LaTeX ソース
\end{lstlisting}

\section{最初の例: テンソルを扱う}

本格的な実装に先立ち、PyTorchテンソルの基本を復習しましょう。テンソルはNumPy配列とほぼ同じように動作しますが、GPU演算と自動微分（autograd）をサポートするという点が異なります。

\begin{lstlisting}[language=Python]
import torch

# テンソルの作成
x = torch.randn(2, 3)         # 標準正規分布からのサンプリング
print(x.shape)                # torch.Size([2, 3])
print(x.dtype)                # torch.float32

# 演算
y = x @ x.T                   # 行列積: [2, 3] @ [3, 2] -> [2, 2]
z = x.sum()                   # すべての元素和（スカラー）

# 自動微分
w = torch.randn(3, requires_grad=True)
loss = (w * w).sum()          # L2 norm^2
loss.backward()               # backward
print(w.grad)                 # d(loss)/d(w) = 2*w
\end{lstlisting}

\begin{infobox}[自動微分（Autograd）]
PyTorchの最も強力な機能の1つ。テンソルで実行されるすべての操作を内部的に「計算グラフ」として記録し、一度にすべてのパラメータの勾配を10回自動的に計算します。私たちは直接微分式を導く必要はありません。この本のすべてのモデルはこのautogradに依存しています。
\end{infobox}

\section{再現性の確保}

ディープラーニング実験はランダム性が多く、再現性が重要である。シードを固定すると、同じコードを何度も実行しても同じ結果が出ます。

\begin{lstlisting}[language=Python]
from makellm.utils import set_seed
set_seed(42)  # Python, NumPy, PyTorch シードを一度に設定
\end{lstlisting}

この関数はPythonの\code{random}、NumPyの\code{numpy.random}、およびPyTorchの\code{torch.manual\_seed}をすべて呼び出します。 CUDAが使用可能な場合は、CUDAシードも設定します。

\section{ユニットテストの実行}

各章のコードは、pytestベースのユニットテストで提供されています。コードを修正したら、必ずテストを実行して回帰がないことを確認しましょう。

\begin{lstlisting}[style=bashstyle]
cd make-llm-basic
pytest tests/ -v
\end{lstlisting}

\begin{lstlisting}[style=bashstyle, basicstyle=\ttfamily\small]
========================= 48 passed in 2.47s =========================
\end{lstlisting}

48のテストがすべて合格すれば環境設定が完了したのだ。

\section{要約}

\begin{itemize}
\item PyTorch 2.xとpytestをインストールしました。
\item リポジトリ構造を理解した:\code{src/makellm/}にモジュールごとにコードが分離されている。
\item テンソル演算とautograd의基本を復習した。
\item 再現性のために\code{set\_seed()}を使用してください。
\item \code{pytest tests/}でユニットテストを実行できます。
\end{itemize}

次の章では、本格的にトルクナイザーを実装します。テキストを数字に置き換えるこの最初の変換は、すべてのLLMの出発点です。
